\documentclass[11pt,letterpaper]{article}
\usepackage[T1]{fontenc}
\usepackage[utf8]{inputenc}
\usepackage{lmodern}
\usepackage{microtype}
\usepackage[margin=1in,headheight=15pt]{geometry}
\usepackage{amsmath,amssymb,amsthm,mathtools,mathrsfs}
\usepackage{enumitem,booktabs,array}
\usepackage[titles]{tocloft}
\setlength{\cftbeforesecskip}{4pt}
\usepackage{xcolor}
\usepackage[colorlinks=true,linkcolor=blue!45!black,citecolor=blue!45!black,urlcolor=blue!45!black]{hyperref}
\hypersetup{
 pdftitle={Exact Hahn Normal Forms in Surcomplex Analysis},
 pdfauthor={Research manuscript prepared with ChatGPT},
 pdfsubject={Integrability, sharp small-divisor thresholds, and support-controlled domains},
 pdfkeywords={surcomplex, Hahn series, Hamiltonian, normal form, small divisors}
}
\usepackage{bookmark}
\usepackage{fancyhdr}
\pagestyle{fancy}
\fancyhf{}
\fancyhead[L]{\small Exact Hahn normal forms}
\fancyhead[R]{\small Surcomplex analysis}
\fancyfoot[C]{\thepage}
\setlength{\emergencystretch}{2em}
\setlist{itemsep=3pt,topsep=5pt}
\numberwithin{equation}{section}
\newtheorem{theorem}{Theorem}[section]
\newtheorem{proposition}[theorem]{Proposition}
\newtheorem{lemma}[theorem]{Lemma}
\newtheorem{corollary}[theorem]{Corollary}
\theoremstyle{definition}
\newtheorem{definition}[theorem]{Definition}
\newtheorem{example}[theorem]{Example}
\theoremstyle{remark}
\newtheorem{remark}[theorem]{Remark}
\newcommand{\C}{\mathbb C}
\newcommand{\R}{\mathbb R}
\newcommand{\N}{\mathbb N}
\newcommand{\Z}{\mathbb Z}
\newcommand{\Q}{\mathbb Q}
\newcommand{\No}{\mathbf{No}}
\newcommand{\E}{\mathscr E}
\newcommand{\Pp}{\mathscr P}
\newcommand{\OO}{\mathcal O}
\newcommand{\mm}{\mathfrak m}
\newcommand{\supp}{\operatorname{supp}}
\newcommand{\ad}{\operatorname{ad}}
\newcommand{\id}{\operatorname{id}}
\newcommand{\nr}{\mathrm{nr}}
\newcommand{\ord}{\operatorname{ord}}
\newcommand{\st}{\operatorname{st}}
\newcommand{\HH}[1]{#1((t^\Gamma))}
\newcommand{\pos}{_{>0}}
\newcommand{\res}{\operatorname{res}}
\newcommand{\wt}{\operatorname{wt}}
\newcommand{\lab}{\mathrm{lab}}
\DeclareMathOperator{\Cent}{Cent}
\title{\textbf{Exact Hahn Normal Forms\newline in Surcomplex Analysis}\\[7pt]
\large Integrability, sharp small-divisor thresholds,\\
and support-controlled domains}
\author{Research manuscript prepared with ChatGPT}
\date{21 September 2026}

\begin{document}
\maketitle
\thispagestyle{empty}
\begin{abstract}
A formal normal form need not converge as an ordinary analytic change of
coordinates. We identify a different, genuinely realizable situation in
surcomplex analysis. Let a nonresonant quadratic Hamiltonian be perturbed by
a positive Hahn series whose coefficients are polynomials in the phase
variables. For arbitrary set-sized ordered value groups and arbitrary
well-ordered positive supports, we construct a canonical Hahn-summable
symplectic normalization. It is an actual invertible analytic change of
coordinates on the entire finite surcomplex phase space, with no arithmetic
bound on the nonzero homological divisors. The resulting commuting actions
generate the whole Poisson centralizer in the entire-coefficient Hahn
algebra.

For entire, rather than polynomial, coefficient functions, we prove an exact
universal criterion: the reciprocal homological divisors must grow at most
exponentially in phase degree. If this fails, a single positive Hahn layer
with an entire coefficient already obstructs every first-order symplectic
normalization with ordinary analytic coefficient germs. An explicit super-Liouville example makes the
obstruction concrete. Finally, a monomial support-and-degree certificate
extends polynomial normalizations to valuation polydisks of possibly
infinite radius; examples show separately why strict positivity,
well-ordering, and finite fibers cannot simply be omitted.

The classical formal normalization algorithm, formal Lie exponentials, and
rank-one non-Archimedean linearization are not claimed as new. The proposed
contributions are the precise realization, threshold, centralizer, and
rescaling statements proved here. Priority has not been established by an
exhaustive literature review. Accompanying code passes 119 exact finite
symbolic checks; these checks are not a substitute for the proofs.
\end{abstract}
\vfill
\noindent\textbf{Keywords.} Surcomplex numbers; Hahn series; Hamiltonian normal
forms; small divisors; exact integrability; symplectic substitutions.
\newpage
{\small\tableofcontents}
\newpage

\section{Introduction and the questions answered}
\label{sec:introduction}

\subsection{Why dynamics is a useful next subject}
The repository motivating this work contains a substantial collection on
surreal and surcomplex analysis. Its reviewed catalogue distinguishes
coherent, common-domain coefficient families from unrelated analytic germs,
and lists treatments of residues, finite deformations, analytic geometry,
divisors, polynomial algebra, and trigonometry \cite{repository}. It does not
list a dedicated dynamics report. We take its coefficient-domain distinction
seriously and investigate a different question: when does a normal form
become an \emph{actual} surcomplex change of variables, rather than merely a
formal expression in the phase variables?

The question is not answered by the algebraic closure of $\No[i]$.
Normalization involves infinitely many substitutions, divisions, and Lie
brackets, so an algebraically closed coefficient field does not by itself
provide a meaningful sum. Nor is a proof that works for a single formal
parameter automatically valid for an arbitrary-rank value group. An argument
that silently assumes $n\eta\to+\infty$ can fail even when all its individual
algebraic identities are correct.

This article answers the following questions in explicitly specified
function classes.
\begin{enumerate}[label=\textbf{Q\arabic*.},leftmargin=3em]
\item Can an infinitesimal Hamiltonian perturbation with arbitrarily many
Hahn scales be normalized and integrated on the whole finite surcomplex
phase space, without a Diophantine condition?
\item Does the answer change when one Hahn coefficient is allowed to be an
entire function with infinitely many Taylor coefficients? Is there a sharp
universal arithmetic threshold?
\item Can the normalizing map be extended to phase points of infinite
magnitude, and can such an extension be certified from the input support
and phase degrees?
\end{enumerate}
The answers are affirmative for Q1 in the polynomial-coefficient category;
Q2 has the exact exponential-divisor criterion of
Theorem~\ref{thm:universal}; and Q3 has the support certificate of
Theorem~\ref{thm:rescale}, with sharp failure examples.

\subsection{Main conclusions in a compact form}
Write $I_j=q_jp_j$ and fix
\begin{equation}\label{eq:H0intro}
 H_0=\sum_{j=1}^d\omega_j I_j,
 \qquad \omega\cdot k\ne0\quad(0\ne k\in\Z^d).
\end{equation}
Here the $\omega_j$ are ordinary complex constants. Suppose
\[
 H=H_0+\sum_{s\in S}t^s R_s(q,p),
 \qquad S\subseteq\Gamma_{>0}\text{ well ordered},
 \qquad \ord_0 R_s\ge3.
\]
Our principal conclusions are as follows.

\paragraph{Polynomial layers.}
If every $R_s$ is a polynomial, no estimate on the sizes of the nonzero
numbers $\omega\cdot k$ is required. There is a uniquely gauge-normalized
positive generator $\chi$ such that
\[
 e^{\ad_\chi}H=N(I),\qquad
 N(I)=\omega\cdot I+\text{positive Hahn terms of action degree at least two}.
\]
Both the resulting map and its inverse are defined at every finite
surcomplex phase point. Their coefficients are still polynomials, although
the degrees need not be uniformly bounded.

\paragraph{Entire layers.}
Define the divisor growth invariant
\begin{equation}\label{eq:Lambdaintro}
 \Lambda(\omega)=\limsup_{D\to\infty}
 \left[\max\left(1,
 \max_{\substack{|\alpha|+|\beta|=D\\\alpha\ne\beta}}
 |\omega\cdot(\alpha-\beta)|^{-1}\right)\right]^{1/D}.
\end{equation}
Universal normalization with entire Hahn coefficients holds exactly when
$\Lambda(\omega)<\infty$. If $\Lambda(\omega)=\infty$, an entire perturbation
$H_0+t^\eta F$ can fail to admit even an analytic first-order normalizing
coefficient. Thus positivity of the Hahn scale cannot cure every
small-divisor problem.

\paragraph{All coherent first integrals.}
For every Hamiltonian successfully normalized above, the functions
$J_j=e^{-\ad_\chi}I_j$ commute, and
\[
 \{F,H\}=0
 \quad\Longleftrightarrow\quad
 F=G(J_1,\ldots,J_d)
\]
for a unique entire-coefficient Hahn series $G$. This is a classification in
a specified algebra, not a claim about arbitrary locally defined functions.

\paragraph{Larger domains.}
For a monomial $t^\gamma q^\alpha p^\beta$ in the perturbation, put
\[
 w_\rho=\gamma+(|\alpha|+|\beta|-2)\rho.
\]
If the shifted monomial family is positive, well ordered, and finite-to-one
at each shifted weight, the same normalization is realizable on
$v(q_j),v(p_j)\ge\rho$. Negative $\rho$ permits infinite phase magnitudes.
The criterion is not merely a lower bound on individual valuations.

\subsection{Relation to known work and the status of novelty}
Classical Hamiltonian normal-form theory already supplies formal
normalizations and formal integrability at a nonresonant quadratic part.
Their ordinary analytic convergence is a separate and difficult issue;
P\'erez-Marco and Krikorian treat convergence/divergence questions in that
classical setting \cite{perezmarco,krikorian}. We neither rebrand the formal
algorithm as a new theorem nor infer ordinary analytic convergence from a
Hahn identity.

Non-Archimedean linearization in equal characteristic zero is also
established. Lindahl's work includes complete real-valued ultrametric fields
and explains why residue-field nonresonance removes the relevant divisor
loss there \cite{lindahl}. Our arbitrary-rank support arguments and the
polynomial-versus-entire distinction should not be confused with discovery
of that rank-one phenomenon.

Exponentiating suitable derivations in algebras with formal sums has a
substantial general theory. Bagayoko, Krapp, Kuhlmann, Panazzolo, and Serra
prove an exponential correspondence in a broad summability setting
\cite{bagayoko}. We give direct proofs for the particular support-controlled
operators needed here, rather than claim a new exponential correspondence.

The exact package of statements developed here---global finite-halo
realization with arbitrary positive Hahn support, the universal
entire-coefficient threshold and its counterexample, the entire-coherent
Poisson centralizer, and compatible monomial-weight domain certificates---was
not located in the consulted primary sources. These are therefore
\emph{proposed contributions with proofs}, not certified bibliographic
firsts. No named, previously published open conjecture is claimed solved.

\section{Hahn algebras and actual surcomplex evaluation}
\label{sec:hahn}

\subsection{Set-sized workspaces}
Fix a nontrivial set-sized ordered abelian group $\Gamma$. Put
\[
 K_\Gamma=\C((t^\Gamma)).
\]
An element is a formal sum $a=\sum_\gamma a_\gamma t^\gamma$ whose nonzero
support is a well-ordered subset of $\Gamma$. Its valuation is
$v(a)=\min\supp(a)$ for $a\ne0$, with $v(0)=+\infty$. Write
\[
 \OO_\Gamma=\{a:v(a)\ge0\},\qquad
 \mm_\Gamma=\{a:v(a)>0\}.
\]
Every element of $\OO_\Gamma$ has a unique decomposition $a=c+u$ with
$c\in\C$ and $u\in\mm_\Gamma$.

For the surcomplex application take $\Gamma$ to be a set-sized ordered
subgroup of $\No$, and interpret $t^\gamma=\omega^{-\gamma}$. The classical
normal-form description of surreal numbers identifies such Hahn series with
surreal real and imaginary parts; a recent treatment records this framework
and its class-size conventions \cite{kaplan}. Divisibility of $\Gamma$ may be
imposed when an algebraically closed workspace is desired, but the proofs
below do not need algebraic closure.

A fixed set of input series and a finite tuple of arguments use only a set
of exponents. Enlarging the workspace to contain those exponents, and if
necessary their divisible hull, is legitimate. All constructions below
commute with such inclusions. Consequently ``every finite surcomplex phase
point'' is read pointwise through a suitable workspace, not by forming a set
of all surcomplex numbers.

\subsection{Two coefficient algebras, not one}
Let
\[
 \Pp_n=\C[z_1,\ldots,z_n],\qquad
 \E_n=\OO(\C^n).
\]
For either $A_n=\Pp_n$ or $A_n=\E_n$, define
\begin{equation}\label{eq:coeffalgebra}
 \HH{A_n}=
 \left\{\sum_{\gamma\in A}t^\gamma f_\gamma(z):
 A\subseteq\Gamma\text{ well ordered},\ f_\gamma\in A_n\right\}.
\end{equation}
The support is understood after zero coefficients are removed. A subscript
$>0$ means that this support is contained in $\Gamma_{>0}$. Differentiation
always acts on the phase variables and fixes the Hahn monomials $t^\gamma$.

In $\HH{\E_n}$ all coefficients have the common domain $\C^n$. We impose
\emph{no} uniform norm bound as $\gamma$ varies. Likewise, in
$\HH{\Pp_n}$ we impose no uniform degree bound. An element of
$\HH{\C\{z\}}$, whose coefficients are unrelated convergent germs, is a
different object. Results about actual evaluation on the finite halo are not
silently transferred to that germ algebra.

\begin{definition}[Strong summability]
A set-indexed family $(F_i)_{i\in J}$ of Hahn series with coefficients in a
fixed vector space is strongly summable if the union of its supports is well
ordered and, for each $\gamma$, only finitely many members have a nonzero
coefficient at $\gamma$. Its sum is then defined coefficientwise by finite
addition.
\end{definition}
Ordinary convergent Taylor series inside a single entire coefficient are
part of ordinary complex analysis. Summation of different Hahn layers is
instead governed by this definition. The two operations must not be
interchanged without a support argument.

\begin{lemma}[Positive-support lemma]\label{lem:neumann}
Let $S\subseteq\Gamma_{>0}$ be well ordered and let
\[
 M(S)=\{0\}\cup\bigcup_{n\ge1}\{s_1+\cdots+s_n:s_j\in S\}.
\]
Then $M(S)$ is well ordered. Each $\gamma\in\Gamma$ is the sum of only
finitely many finite ordered words in $S$. If $A\subseteq\Gamma$ is well
ordered, then $A+M(S)$ is well ordered and, for fixed $\gamma$, there are only
finitely many pairs consisting of $a\in A$ and a word in $S$ with total
$a+s_1+\cdots+s_n=\gamma$.
\end{lemma}
This is the standard positive-support form of Neumann's lemma
\cite{neumann}. The last assertion follows from the finite-decomposition
property for sums of two well-ordered subsets and the word assertion. It
applies to arbitrary ordered groups, not just subgroups of $\R$.

\begin{remark}[No hidden rank-one hypothesis]\label{rem:rank}
In $\Gamma=\Z^2$ with lexicographic order, put $\eta=(0,1)$ and
$\theta=(1,0)$. Then $n\eta<\theta$ for every ordinary integer $n$.
Nevertheless $\sum_{n\ge0}t^{n\eta}$ is a valid Hahn series, and its product
with $1-t^\eta$ is $1$. There are infinitely many terms below the cutoff
$\theta$, so ``only finitely many terms below each cutoff'' is not our
summability criterion. Likewise, the partial sums need not converge in the
valuation topology. All arguments below use finite contributions at a
\emph{fixed exponent}.
\end{remark}

\subsection{Evaluation on the finite halo}
For $F=\sum_\gamma t^\gamma f_\gamma\in\HH{\E_n}$ and
$z=a+u\in\OO_\Gamma^n$, where $a\in\C^n$ and $u\in\mm_\Gamma^n$, define
\begin{equation}\label{eq:eval}
 F^\#(a+u)=
 \sum_{\gamma,\alpha\in\N^n}
 t^\gamma\frac{\partial^\alpha f_\gamma(a)}{\alpha!}u^\alpha.
\end{equation}

\begin{proposition}[Faithful analytic realization]\label{prop:eval}
Expression \eqref{eq:eval} is strongly summable. It defines a faithful
$K_\Gamma$-algebra homomorphism from $\HH{\E_n}$ to functions on
$\OO_\Gamma^n$. It respects phase derivatives and substitutions by
positive near-identity Hahn maps. At every finite point its infinitesimal
Taylor expansions have the expected coefficients.
\end{proposition}
\begin{proof}
Let $A$ be the support of $F$ and let $T$ be the union of the supports of the
nonzero components of $u$. This is a finite union of well-ordered positive
sets and hence is well ordered and positive. Every monomial contribution to
\eqref{eq:eval} has exponent in $A+M(T)$. For a fixed exponent,
Lemma~\ref{lem:neumann} allows only finitely many choices of $\gamma$ and a
word of exponents from $T$. The word length bounds $|\alpha|$, and there are
only finitely many coordinate assignments of a fixed finite word. Thus the
family in \eqref{eq:eval}, with its products expanded, is strongly summable.

The ordinary Taylor identities for products and derivatives now apply, since
every regrouping at a Hahn exponent is finite. The same reasoning, adding
the positive support of a near-identity substitution to $T$, proves the
composition assertion from the ordinary Taylor chain rule.

For faithfulness, take the least Hahn exponent $\gamma_0$ of a nonzero $F$.
The nonzero entire function $f_{\gamma_0}$ is nonzero at some ordinary point
$a\in\C^n$. At that point the coefficient of $t^{\gamma_0}$ in $F^\#(a)$ is
nonzero.

Finally, at a finite point $z$ and an infinitesimal increment $h$, the same
regrouping proves the Taylor formula with coefficients
$(\partial^\alpha F)^\#(z)/\alpha!$. For example, if $a_0=\min\supp F$ and
$\mu=\min_j v(h_j)>0$, its terms of total increment degree at least two have
valuation at least $a_0+2\mu$. This also yields the usual local derivative
estimate after division by a scalar increment. No limit of Hahn partial
sums has been used.
\end{proof}

\begin{remark}
The word ``analytic'' in this article refers to this entire-coefficient,
Taylor-compatible realization. It does not assert that an ordinary real
interval maps continuously and nontrivially into the fine topology of the
whole proper-class field $\No[i]$. Nor does it identify phase differentiation
with a derivation of the coefficient field of surreal numbers.
\end{remark}

\section{Support-controlled substitutions and Lie maps}
\label{sec:lie}

\subsection{Inversion without a convergence assumption}
\begin{proposition}[Positive near-identity maps]\label{prop:inverse}
Let $U\in\HH{\E_n}_{>0}^{\,n}$. Then $\Phi(z)=z+U(z)$ has a unique inverse
of the form $z+V(z)$ with $V\in\HH{\E_n}_{>0}^{\,n}$. Both are bijections
of $\OO_\Gamma^n$, preserving standard parts. If $U$ has polynomial
coefficients, so does $V$.
\end{proposition}
\begin{proof}
Let $S$ be the union of the supports of the components of $U$ and put
$M=M(S)$. The equation for the inverse is
\begin{equation}\label{eq:inverse}
 V(z)=-U(z+V(z)).
\end{equation}
Expand the right side by Taylor's formula. Its term independent of $V$ is
$-U(z)$, and every other term has at least one positive input from $U$ and at
least one from $V$. At exponent $\gamma>0$, all occurrences of coefficients
of $V$ therefore have exponents strictly smaller than $\gamma$. Recursion on
the well-ordered set $M$ uniquely determines $V_\gamma$.

At a fixed exponent, only finitely many Taylor terms and input words can
occur by Lemma~\ref{lem:neumann}. Their coefficients use only derivatives,
products, and finite sums of the coefficients of $U$. They are entire, or
polynomial in the polynomial case. This also proves existence at limit
positions of $M$: no topological limit is requested there.

Equation \eqref{eq:inverse} gives a right inverse. The same construction for
a left inverse, followed by associativity of supported substitution, shows
that the two are equal. Alternatively their composition equations and the
least-exponent uniqueness argument give the same conclusion. Realization
by Proposition~\ref{prop:eval} proves the asserted actual bijections.
\end{proof}

\subsection{Poisson brackets and their exponentials}
In $2d$ phase variables $q_1,\ldots,q_d,p_1,\ldots,p_d$, use
\begin{equation}\label{eq:pb}
 \{F,G\}=\sum_{j=1}^d
 (F_{q_j}G_{p_j}-F_{p_j}G_{q_j}),\qquad
 \ad_\chi F=\{\chi,F\}.
\end{equation}
This sign convention will be maintained throughout.

\begin{proposition}[Supported Hamiltonian exponential]\label{prop:exp}
If $\chi\in\HH{\E_{2d}}_{>0}$, then for every
$F\in\HH{\E_{2d}}$ the family
\[
 \left(\frac{\ad_\chi^nF}{n!}\right)_{n\ge0}
\]
is strongly summable. Its sum $T_\chi F=e^{\ad_\chi}F$ is an algebra and
Poisson automorphism, with inverse $T_{-\chi}$. If
\[
 \Phi=(T_\chi q_1,\ldots,T_\chi q_d,
             T_\chi p_1,\ldots,T_\chi p_d),
\]
then $T_\chi F=F\circ\Phi$. The realized map $\Phi$ is a symplectic
bijection of the finite halo. Polynomial coefficients are preserved.
\end{proposition}
\begin{proof}
If $A=\supp F$ and $S=\supp\chi$, each term of $\ad_\chi^nF$ has support
in $A+nS$. A fixed exponent receives contributions for only finitely many
words in $S$ attached to an exponent in $A$. Differentiation does not change
Hahn weights. Thus Lemma~\ref{lem:neumann} proves strong summability,
including all the regroupings used next.

The Leibniz rule gives
\[
 \ad_\chi^n(FG)=\sum_{j=0}^n\binom nj
       (\ad_\chi^jF)(\ad_\chi^{n-j}G).
\]
After division by $n!$ and supported summation, this is multiplicativity of
$T_\chi$. The Jacobi identity gives the identical binomial rule for the
Poisson bracket, so $T_\chi\{F,G\}=\{T_\chi F,T_\chi G\}$.
The binomial identity for $\exp(D)\exp(-D)$, justified by the same support
argument, gives the inverse.

For an ordinary entire function, the formal chain rule for the derivation
$\ad_\chi$ identifies $T_\chi F$ with the Taylor substitution by the
coordinate tuple $\Phi$. One may check this at each power of an auxiliary
formal scalar multiplying $\chi$, where it is the usual formal flow
identity, and then evaluate that scalar at $1$ using the proved word
summability. Applying this coefficientwise gives the assertion for general
$F$. The coordinate Poisson relations are preserved, proving
symplecticity. Finally $\Phi=\id+\text{positive terms}$, so
Propositions~\ref{prop:inverse} and \ref{prop:eval} give the actual maps.
\end{proof}

\subsection{The finite-expression certificate}
We will repeatedly use a little more than the existence of a sum. Label an
input perturbation coefficient by its positive Hahn exponent. A term in a
normalizing coefficient is built by a finite expression tree whose leaves
are labeled inputs and whose internal operations are Poisson brackets,
linear projections, scalar division, or fixed linear solution operators.
Its weight is the sum of its leaf weights.

For a fixed total weight there are finitely many leaf words. In the
recursion below, every genuinely nonlinear internal step has at least two
positive inputs; linear operators only decorate these steps or the leaves.
Thus a fixed leaf word admits only finitely many relevant tree shapes.
This is a finite certificate for every coefficient, even when the set of
smaller weights is infinite. It will also be the reason that a later
reordering of weights under rescaling is legitimate.

\section{The homological operator and a sharp divisor threshold}
\label{sec:divisors}

\subsection{Resonances and the action projection}
Fix the nonresonant frequency vector of \eqref{eq:H0intro}. Define
\[
 L F=\{F,H_0\}.
\]
A direct calculation gives
\begin{equation}\label{eq:homological}
 L(q^\alpha p^\beta)
   =\omega\cdot(\alpha-\beta)\,q^\alpha p^\beta.
\end{equation}
Nonresonance means that its kernel consists exactly of the monomials with
$\alpha=\beta$.

For a Taylor expansion at the origin, define
\begin{equation}\label{eq:projection}
 \Pi\left(\sum_{\alpha,\beta}c_{\alpha\beta}q^\alpha p^\beta\right)
  =\sum_\alpha c_{\alpha\alpha}I^\alpha,
 \qquad Q=1-\Pi.
\end{equation}
The letter $Q$ denotes this projection only; in examples a coordinate named
$Q_1$ will always have a subscript or be explicitly introduced.

\begin{lemma}[Entire action projection]\label{lem:projection}
Projection $\Pi$ maps $\E_{2d}$ to functions $g(I)$ with $g\in\E_d$, and
$g$ is unique. It preserves polynomial coefficients, does not decrease
phase vanishing order, and commutes with $L$. On $Q\Pp_{2d}$ the inverse
$L^{-1}$ is always defined and polynomial, without an arithmetic estimate.
\end{lemma}
\begin{proof}
For an entire $f$, its Taylor coefficients satisfy for every $R>0$
\[
 |c_{\alpha\beta}|\le C_R R^{-|\alpha|-|\beta|}
\]
by Cauchy's estimate on a phase polydisk of radius $R$. In particular
$|c_{\alpha\alpha}|\le C_R(R^2)^{-|\alpha|}$. Since $R$ is arbitrary, the
action series is entire. Distinct action monomials are distinct phase
monomials, proving uniqueness. The remaining projection assertions follow
termwise. For a polynomial the inverse divides finitely many nonresonant
coefficients by nonzero complex constants, so it is a polynomial.
\end{proof}
One can also realize $\Pi$ by averaging over the ordinary torus action
$q_j\mapsto e^{i\theta_j}q_j$, $p_j\mapsto e^{-i\theta_j}p_j$. No
surcomplex contour integration is needed for this observation.

\subsection{The exact analytic condition}
For $D\ge1$, let
\begin{equation}\label{eq:MD}
 M_D(\omega)=\max\left(1,
 \max_{\substack{|\alpha|+|\beta|=D\\\alpha\ne\beta}}
      |\omega\cdot(\alpha-\beta)|^{-1}\right),
 \qquad \Lambda(\omega)=\limsup_D M_D(\omega)^{1/D}.
\end{equation}
All inner maxima are finite. Finiteness of $\Lambda$ is equivalent to the
existence of constants $C,A\ge1$ with $M_D\le CA^D$ for every $D$.

\begin{theorem}[Sharp homological criterion]\label{thm:homological}
For a nonresonant $\omega$, the following are equivalent:
\begin{enumerate}[label=(\roman*)]
\item $\Lambda(\omega)<\infty$.
\item The coefficientwise operator $L^{-1}$ maps every nonresonant entire
function to an entire function.
\item The coefficientwise operator $L^{-1}$ maps every nonresonant
convergent germ at the origin to a convergent germ, with a possible loss of
radius.
\end{enumerate}
If these fail, there is a nonresonant entire function $F$ of order at least
three for which $L^{-1}F$ has no positive radius of convergence.
\end{theorem}
\begin{proof}
Suppose $M_D\le CA^D$. If the coefficients of $F$ obey a Cauchy estimate
$|f_{\alpha\beta}|\le C_RR^{-D}$, those of $L^{-1}F$ obey
\[
 |(L^{-1}F)_{\alpha\beta}|\le CC_R(A/R)^D.
\]
The number of monomials of total degree $D$ in $2d$ variables is
$\binom{D+2d-1}{2d-1}$, which grows polynomially in $D$. Hence this bound
proves absolute convergence on every polydisk of radius less than $R/A$.
For a germ choose one admissible $R$; for an entire function let $R$ be
arbitrarily large. This proves (i)$\Rightarrow$(ii),(iii).

Conversely, suppose $\Lambda=\infty$. Select distinct degrees $D_j\to\infty$,
with $D_j\ge\max(j,3)$, and nonresonant monomials
$m_j=q^{\alpha_j}p^{\beta_j}$ of those degrees such that
\[
 \delta_j=|\omega\cdot(\alpha_j-\beta_j)|,
 \qquad \delta_j^{-1/D_j}\ge j^4.
\]
Discarding an initial term if necessary causes no problem. Define
\begin{equation}\label{eq:sparseF}
 F=\sum_j \delta_j^{1/2}m_j.
\end{equation}
On a polydisk of radius $R$, the absolute size of the $j$th term is at most
$(R/j^2)^{D_j}$, and the tail is bounded by a geometric series because
$D_j\ge j$. Thus $F$ is entire. In the putative inverse the absolute
coefficient of $m_j$ is $\delta_j^{-1/2}$, whose $D_j$th root is at least
$j^2$. These coefficients violate every Cauchy bound on a positive-radius
polydisk. The inverse is not a convergent germ. Since $F$ is itself a germ,
both (ii) and (iii) fail. Its order and nonresonance were built into the
selection.
\end{proof}

\begin{remark}[What condition is, and is not, being measured]
A polynomial lower bound on $|\omega\cdot k|$ implies $\Lambda=1$, but is
far stronger than necessary here. Finite exponential growth is the exact
condition for this \emph{single coefficientwise homological operator} to
preserve all entire functions. The nonlinear Hahn theorem below repeatedly
uses that operator, but it requires no uniform norm estimate across Hahn
layers. This is why the criterion must not be advertised as an ordinary
complex-analytic convergence criterion of Bruno or KAM type.
\end{remark}

\section{Arbitrary-support Hamiltonian normalization}
\label{sec:normalform}

\subsection{The main existence and realization theorem}
\begin{theorem}[Exact positive-Hahn normal form]\label{thm:normalform}
Let $\Gamma$ be as above, let $\omega$ satisfy \eqref{eq:H0intro}, and let
\begin{equation}\label{eq:Hinput}
 H=H_0+R,\qquad
 R=\sum_{s\in S}t^s R_s(q,p),\qquad
 S\subseteq\Gamma_{>0}\text{ well ordered},\quad \ord_0 R_s\ge3.
\end{equation}
Assume either of the following hypotheses:
\begin{enumerate}[label=(\alph*)]
\item every $R_s$ is a polynomial, with no further restriction on $\omega$;
\item every $R_s$ is entire and $\Lambda(\omega)<\infty$.
\end{enumerate}
There is a unique positive generator $\chi$ in the corresponding coefficient
algebra satisfying
\begin{equation}\label{eq:gauge}
 \Pi\chi=0,\qquad \ord_0\chi_\gamma\ge3,
 \qquad e^{\ad_\chi}H=H_0+B(I),
\end{equation}
where $B$ has positive Hahn support and its coefficients have action
vanishing order at least two. The supports of $\chi$ and $B$ are contained
in $M(S)\setminus\{0\}$.

The coordinate map $\Phi=e^{\ad_\chi}(q,p)$ and its inverse are symplectic
bijections of $\OO_\Gamma^{2d}$. They fix the origin, are tangent to the
identity there, preserve standard parts everywhere, and satisfy the actual
identity
\begin{equation}\label{eq:actualnormal}
 H^\#\circ\Phi^\#=(H_0+B(I))^\#.
\end{equation}
All statements commute with enlargement of the Hahn workspace.
\end{theorem}

\begin{proof}
We first construct the generator, keeping summability separate from the
normal-form calculation. Expanding the Lie exponential yields
\begin{equation}\label{eq:normalexpansion}
 e^{\ad_\chi}(H_0+R)=H_0+R+L\chi+\mathcal B(\chi,R),
\end{equation}
where
\begin{equation}\label{eq:Bdef}
 \mathcal B(\chi,R)=
 \sum_{n\ge1}\frac{\ad_\chi^n R}{n!}
 +\sum_{n\ge2}\frac{\ad_\chi^n H_0}{n!}.
\end{equation}
Every term in $\mathcal B$ involves at least two positive factors, counting
occurrences of $R$ and $\chi$. Since $\Pi L=0$ and $\Pi\chi=0$, the
nonresonant normal-form equation is exactly
\begin{equation}\label{eq:chifixed}
 \chi=-L^{-1}Q\bigl(R+\mathcal B(\chi,R)\bigr).
\end{equation}

Put $M=M(S)$. At a weight $\gamma\in M\setminus\{0\}$, the coefficient of
$\mathcal B$ depends only on $\chi_\mu$ with $0<\mu<\gamma$: an occurrence
of $\chi_\gamma$ would be accompanied by another positive factor and hence
would have total weight greater than $\gamma$. Thus well-founded recursion
on $M$ determines
\begin{equation}\label{eq:coefficientrecursion}
 \chi_\gamma=-L^{-1}Q\left(R_\gamma+
      [t^\gamma]\mathcal B(\chi_{<\gamma},R)\right).
\end{equation}
This formula remains meaningful when there are infinitely many predecessors
of $\gamma$. Lemma~\ref{lem:neumann} says that only finitely many of them
can participate in a word of total weight $\gamma$.

More explicitly, each coefficient obtained by recursion is a finite sum of
labeled expression trees described in Section~\ref{sec:lie}. Their leaves
are input coefficients $R_s$, their total weight is $\gamma$, and their
internal nonlinear nodes have at least two positive inputs. The finitely
many words of total weight $\gamma$ bound the number of leaves and their
choices. There are finitely many tree shapes and bracket orderings for
those leaves. Consequently every coefficient uses finitely many
applications of derivatives, multiplication, $Q$, and $L^{-1}$. This proves
both the claimed support bound and the absence of an unproved infinite
operation within a Hahn coefficient.

In case (a), these operations preserve polynomials by
Lemma~\ref{lem:projection}. In case (b), they preserve entire functions by
Theorem~\ref{thm:homological}. The generator therefore belongs to the
asserted algebra. Every application of $L^{-1}Q$ gives $\Pi\chi_\gamma=0$.
The elementary estimate
\[
 \ord_0\{F,G\}\ge\ord_0F+\ord_0G-2
\]
shows inductively that each $\chi_\gamma$ has order at least three. Terms in
$\mathcal B$ have order at least four. The resonant part
\begin{equation}\label{eq:Bnormal}
 B=\Pi\bigl(R+\mathcal B(\chi,R)\bigr)
\end{equation}
thus has phase order at least four, since a resonant monomial has even
phase degree and $R$ has order at least three. By
Lemma~\ref{lem:projection}, it is an entire or polynomial function of the
actions at every Hahn weight, of action order at least two.

Equation \eqref{eq:chifixed} removes the nonresonant part of
\eqref{eq:normalexpansion}; hence \eqref{eq:gauge} follows as a Hahn
identity. The exponential is defined by Proposition~\ref{prop:exp}.

For uniqueness, take two positive gauge-normalized generators and let
$\gamma$ be the least exponent of their difference. In the difference of
their equations \eqref{eq:chifixed}, all terms coming from $\mathcal B$
have weight strictly greater than $\gamma$, because they involve that
difference and another positive factor. At weight $\gamma$ we therefore
obtain $L(\chi_\gamma-\widetilde\chi_\gamma)=0$. The difference is
nonresonant, on which $L$ is injective by \eqref{eq:homological}. This is a
contradiction. This argument also rules out additional support outside
$M(S)$ in a competing generator.

Finally Proposition~\ref{prop:exp} gives the symplectic maps and inverse.
Because $\chi$ has phase order at least three, their nonlinear coordinate
terms have order at least two. Proposition~\ref{prop:eval} turns the
normal-form identity into \eqref{eq:actualnormal}. No construction used any
exponent outside the additive monoid generated by the input support, so the
same finite coefficient formulas work in every enlarged workspace.
\end{proof}

\subsection{What has become exact}
In a classical analytic problem, the generator can be formal in the phase
variables and have zero radius of convergence. In case (a) above, each Hahn
coefficient of the generator is instead a \emph{finite polynomial}; its
coefficients may be enormous ordinary complex numbers, but they are finite
constants at their Hahn layer. An arbitrarily large ordinary number cannot
cancel the positive valuation of $t^\gamma$. Evaluation on the finite halo
is therefore controlled by supports rather than by a numerical norm of the
normalizing coefficients.

For case (b), infinitely many phase monomials occur within a single Hahn
layer. Their ordinary convergence must be settled \emph{before} Hahn
summability can be invoked. The homological threshold does exactly this.
It does not require a common numerical bound on all the resulting entire
coefficients, because the common domain is already all of $\C^{2d}$.

\begin{corollary}[Finite surcomplex phase space]\label{cor:surcomplex}
For polynomial layers, and for entire layers satisfying the threshold,
Theorem~\ref{thm:normalform} realizes the normalization at every finite
point of $\No[i]^{2d}$. The maps obtained in different set-sized workspaces
agree on their common domain.
\end{corollary}
\begin{proof}
Given a phase tuple, enlarge the workspace to contain all of its exponents
and those of the input. Apply Theorem~\ref{thm:normalform}. Its coefficient
formulas and the Taylor evaluation are unchanged by the enlargement.
\end{proof}

\begin{remark}[Uniqueness has a gauge]
The unique object in Theorem~\ref{thm:normalform} is the single exponential
generator with $\Pi\chi=0$, and the normal form produced in this gauge.
Arbitrary symplectic changes of variables normalizing a Hamiltonian need not
be unique. In particular, transformations generated by functions of the
actions can preserve a normal form. We do not suppress this freedom by an
unstated uniqueness claim.
\end{remark}

\section{Exact integrability and the whole Poisson centralizer}
\label{sec:integrability}

\subsection{Commuting actions}
Write $T=e^{\ad_\chi}$ and $N=TH$ for the objects just constructed. Define
\begin{equation}\label{eq:actionsJ}
 J_j=T^{-1}I_j=e^{-\ad_\chi}I_j.
\end{equation}

\begin{theorem}[Exact integrability]\label{thm:integrability}
The functions $J_1,\ldots,J_d$ are entire-coefficient Hahn functions, with
polynomial coefficients in case (a), and satisfy
\begin{equation}\label{eq:integrals}
 \{H,J_j\}=0,\qquad \{J_j,J_k\}=0,\qquad H=N(J_1,\ldots,J_d).
\end{equation}
They are functionally independent wherever, in the normalized coordinates,
none of the pairs $(q_j,p_j)$ is $(0,0)$. These statements hold as realized
identities on the finite halo.
\end{theorem}
\begin{proof}
The functions $I_j$ commute with one another and with every function of the
actions, including $N$. Applying the Poisson automorphism $T^{-1}$ proves
the first two identities. Applying the algebra substitution $T^{-1}$ to
$N(I)$ proves the third. The $d$ differentials
\[
 dI_j=p_j\,dq_j+q_j\,dp_j
\]
have disjoint coordinate pairs, and are linearly independent exactly when
each is nonzero. Pullback by the invertible differential of $\Phi^{-1}$
preserves that rank. Realization respects all these identities by
Proposition~\ref{prop:eval}.
\end{proof}

\subsection{Every coherent first integral is an action function}
It is important that a classification, not merely the construction of $d$
first integrals, is available.

\begin{theorem}[Poisson centralizer]\label{thm:centralizer}
In the algebra $\HH{\E_{2d}}$ one has
\begin{equation}\label{eq:centralizer}
 \Cent(H):=\{F:\{F,H\}=0\}
 =\{G(J_1,\ldots,J_d):G\in\HH{\E_d}\}.
\end{equation}
The representing series $G$ is unique. Its Hahn support need not be
positive. Thus \eqref{eq:centralizer} describes all first integrals that
admit an entire-coefficient Hahn representation on the finite halo.
\end{theorem}
\begin{proof}
It suffices to classify the centralizer of $N(I)=H_0+B(I)$, because $T$ is
a Poisson automorphism. Let $\widetilde F=TF$ and suppose
$\{\widetilde F,N\}=0$. Apply $\Pi$ coefficientwise. By
Lemma~\ref{lem:projection}, $\Pi\widetilde F=G(I)$ for a unique
$G\in\HH{\E_d}$. Every such action function commutes with $N$, so
\[
 U:=\widetilde F-\Pi\widetilde F
\]
is nonresonant coefficientwise and also commutes with $N$.

If $U\ne0$, let $a$ be the least exponent of its support. The coefficient
at $t^a$ in $\{U,B\}$ is zero: all exponents of $B$ are positive, and all
exponents of $U$ are at least $a$. The coefficient at $t^a$ in
$\{U,N\}=0$ is therefore
\[
 L U_a=0.
\]
But $U_a$ is nonresonant and nonzero, contradicting
\eqref{eq:homological}. Hence $U=0$ and $\widetilde F=G(I)$.
Applying $T^{-1}$ gives $F=G(J)$.

Conversely every $G(J)$ commutes with $H$ by the same Poisson substitution.
Its definition is legitimate even for arbitrary well-ordered support of
$G$: write $J=I+\Delta$ with positive $\Delta$, use the ordinary Taylor
expansion of each entire coefficient of $G$ at $I$, and apply
Lemma~\ref{lem:neumann}. Uniqueness follows either by applying $T$ or by the
injectivity of $g(I)\mapsto g$ in Lemma~\ref{lem:projection}. Finally
Proposition~\ref{prop:eval} is faithful, so no additional identity appears
merely by passing to actual finite-halo functions.
\end{proof}

\begin{remark}[The coefficient class is essential]
The theorem does not classify first integrals defined only near a particular
point, multivalued functions, or arbitrary functions that happen to be
locally constant in a fine topology. Nor does it say that all first
integrals have polynomial coefficients. The entire-coefficient class is
large enough to contain arbitrary entire functions of the actions, but
rigid enough for the leading nonresonant coefficient argument to work.
\end{remark}

\subsection{Finite ordinary time}
Let
\[
 \Omega_j(I)=\frac{\partial N}{\partial I_j}
            =\omega_j+\Delta_j(I),\qquad v(\Delta_j)>0
\]
where the last notation refers to positive Hahn support. For an ordinary
complex time $\tau$, define
\begin{align}
 q_j(\tau)&=e^{\tau\omega_j}
       \exp_{\mathrm H}(\tau\Delta_j(I))q_j(0),\label{eq:qflow}\\
 p_j(\tau)&=e^{-\tau\omega_j}
       \exp_{\mathrm H}(-\tau\Delta_j(I))p_j(0),\label{eq:pflow}
\end{align}
where $\exp_{\mathrm H}(U)=\sum_{n\ge0}U^n/n!$ for positive Hahn $U$.

\begin{proposition}[Exact ordinary-time evolution]\label{prop:flow}
Equations \eqref{eq:qflow}--\eqref{eq:pflow} give an exact flow of $N$ on the
finite halo for every $\tau\in\C$. Conjugation by $\Phi$ gives the flow of
$H$. The coefficients depend holomorphically on $\tau$, and the flow law
holds as an identity of Hahn-coherent functions.
\end{proposition}
\begin{proof}
The actions are constant in \eqref{eq:qflow}--\eqref{eq:pflow}, because the
two factors for each $j$ are inverse. Thus $\Delta_j(I)$ is constant along
the displayed curve. Differentiating coefficientwise in ordinary time gives
\[
 \dot q_j=\Omega_j(I)q_j=\frac{\partial N}{\partial p_j},\qquad
 \dot p_j=-\Omega_j(I)p_j=-\frac{\partial N}{\partial q_j}.
\]
The Hahn exponentials are strongly summable by positivity; each coefficient
is an entire function of $\tau$ and the phase variables. Their constant
Hahn terms are nonzero ordinary exponentials, so finite phase points remain
finite. The usual exponential identities, justified coefficientwise by
finite Hahn regroupings, prove the flow law. Symplectic conjugacy transports
the equations to $H$.
\end{proof}
This is not an assertion about continuous ordinary-time paths in the fine
topology of the entire proper-class field. It is a coefficientwise
holomorphic-time identity. No infinite surreal time, global surreal
exponential of an arbitrary surcomplex argument, or intrinsic coefficient
field derivation is needed.

\section{Failure for entire layers: an exact converse}
\label{sec:counterexample}

\subsection{The universal statement}
\begin{theorem}[Sharp universal normalization criterion]\label{thm:universal}
For a fixed nonresonant frequency vector $\omega$, the following are
equivalent:
\begin{enumerate}[label=(\roman*)]
\item $\Lambda(\omega)<\infty$.
\item Every positive Hahn perturbation with entire coefficients of phase
order at least three admits the entire-coefficient gauge-normalized
generator and finite-halo realization of Theorem~\ref{thm:normalform}, for
every set-sized ordered value group.
\item Every one-parameter Hamiltonian $H_0+\varepsilon F$, with $F$ entire
of order at least three, admits a symplectic formal change
$\Psi=\id+O(\varepsilon)$ with convergent-germ coefficients at the origin
that transforms it into a function of the actions.
\end{enumerate}
If (i) fails, failure of (iii) is already forced at the coefficient of
$\varepsilon$. There is also no positive, gauge-normalized entire Hahn
generator for the same input in any workspace containing its scale.
\end{theorem}
\begin{proof}
The implication (i)$\Rightarrow$(ii) is Theorem~\ref{thm:normalform};
(ii)$\Rightarrow$(iii) follows by taking $\Gamma=\Z$ and
$\varepsilon=t$.

Suppose $\Lambda=\infty$ and choose the entire nonresonant $F$ constructed
in Theorem~\ref{thm:homological}. If a symplectic change as in (iii) exists,
write $\Psi(z)=z+\varepsilon V(z)+O(\varepsilon^2)$, with $V$ a
holomorphic vector-field germ. The first-order symplectic condition says
that $V$ is a symplectic vector field. Its contraction with the constant
symplectic form is a closed holomorphic one-form. A closed holomorphic
one-form germ is exact: a primitive is given on a small star-shaped
polydisk by the radial homotopy integral. Thus, with the sign chosen to
match \eqref{eq:pb}, there is a holomorphic germ $\chi_1$ such that
\[
 V(G)=\{\chi_1,G\}
\]
for every holomorphic germ $G$.

The coefficient of $\varepsilon$ in
$(H_0+\varepsilon F)\circ\Psi$ is $F+L\chi_1$. If the transformed
Hamiltonian is a function of the actions, its nonresonant part must vanish:
\[
 L(Q\chi_1)=-F.
\]
The projection $Q$ preserves convergent germs, by the same Cauchy estimate
used for entire functions in Lemma~\ref{lem:projection}. Coefficientwise
injectivity of $L$ on nonresonant series forces
$Q\chi_1=-L^{-1}F$, which is not a convergent germ. This contradiction
proves (iii)$\Rightarrow$(i).

For the final assertion consider a putative positive gauge-normalized
entire Hahn generator for $H_0+t^\eta F$. If it had a least nonzero
coefficient below $\eta$, its normal-form equation at that weight would
be $L\chi_\gamma=0$, impossible in the gauge. Thus it has no such
coefficient. At weight $\eta$ the nonlinear terms have larger weight, and
one again obtains $\chi_\eta=-L^{-1}F$, which is not entire. This argument
does not depend on the rank or divisibility of the enlarged value group.
\end{proof}

\subsection{An explicit super-Liouville Hamiltonian}
The sparse construction in the preceding theorem depends on a choice of
bad divisors. The following example specifies them directly.

Define integers recursively by
\begin{equation}\label{eq:Liouv}
 a_1=1,\qquad a_{n+1}=10^{4a_n},\qquad Q_n=10^{a_n},\qquad
 \alpha=\sum_{n\ge1}10^{-a_n}.
\end{equation}
Let $P_n/Q_n=\sum_{j=1}^n10^{-a_j}$, so $P_n$ is an integer, and put
$D_n=P_n+Q_n$. Since $0<\alpha<1$, we have $D_n<2Q_n$.
The positive tail satisfies
\begin{equation}\label{eq:Liouvbound}
 0<\alpha Q_n-P_n
 <2Q_n10^{-a_{n+1}}
 =2Q_n10^{-Q_n^4}
 <e^{-D_n^2}.
\end{equation}
The first bound follows by dominating the very rapidly decreasing tail by
twice its first term. For the last bound, $Q_n\ge10$ and
$\log(2Q_n)-Q_n^4\log 10<-4Q_n^2<-D_n^2$.

The number $\alpha$ is irrational. Indeed, if $\alpha=A/B$ were rational,
then the positive number $\alpha Q_n-P_n$ would be at least $1/B$ for every
$n$, contradicting \eqref{eq:Liouvbound}. Hence $\omega=(1,\alpha)$ is
nonresonant.

In two degrees of freedom define
\begin{equation}\label{eq:explicitF}
 F(q_1,q_2,p_1,p_2)
   =\sum_{n\ge1}e^{-D_n^2/2}q_1^{P_n}p_2^{Q_n},
 \qquad H=I_1+\alpha I_2+\varepsilon F.
\end{equation}

\begin{proposition}[Entire input, nonanalytic first normalizer]
\label{prop:explicitbad}
The coefficient $F$ in \eqref{eq:explicitF} is entire and has phase order at
least three. The Hamiltonian has no symplectic normalization
$\id+O(\varepsilon)$ with convergent-germ coefficients to an action-only
normal form. Its first nonresonant homological solution has zero radius of
convergence.
\end{proposition}
\begin{proof}
On a phase polydisk of radius $R$, the $n$th monomial is bounded by
$e^{-D_n^2/2}R^{D_n}$. The distinct integers $D_n$ tend to infinity, so these
bounds are summable for every finite $R$. Thus $F$ is entire. Already
$D_1=11$, proving the order assertion.

By \eqref{eq:homological}, the homological divisor of the $n$th monomial is
$P_n-\alpha Q_n$. The absolute coefficient after division is greater than
$e^{D_n^2/2}$ by \eqref{eq:Liouvbound}. Its $D_n$th root tends to infinity,
which excludes every positive analytic radius. The first-order argument of
Theorem~\ref{thm:universal} applies.
\end{proof}

There is an additional useful feature: both $F$ and its formal homological
solution depend only on $q_1$ and $p_2$. Any two functions of those variables
have zero Poisson bracket. Therefore, in the purely formal phase category,
$\chi=-\varepsilon L^{-1}F$ removes the perturbation with no higher Lie
corrections at all. The obstruction is not divergence of a long nonlinear
iteration. It is already the failure of a single homological inverse to
preserve the specified analytic coefficient class.

\subsection{Why polynomial layers escape the example}
If instead each monomial of \eqref{eq:explicitF} is assigned a distinct
increasing positive Hahn weight, for example
\[
 \widehat H=H_0+\sum_{n\ge1}t^n
 e^{-D_n^2/2}q_1^{P_n}p_2^{Q_n},
\]
then each Hahn layer is a polynomial and Theorem~\ref{thm:normalform}(a)
applies, despite the same frequency vector. Dividing an individual
coefficient by its tiny ordinary divisor still gives an ordinary complex
number, not a failure of a phase Taylor series within that layer.

This is not a contradiction: $\widehat H$ is a different Hamiltonian, with
a different distribution of monomials among the infinitesimal scales. It
shows exactly where Hahn support and ordinary analytic convergence interact.

The divergent formal generator in Proposition~\ref{prop:explicitbad}
can still be evaluated at infinitesimal phase points: its ordinary complex
Taylor coefficients are summed against positive phase supports by
Lemma~\ref{lem:neumann}. The counterexample therefore does not exclude
normalization on an infinitesimal monad. Its obstruction is to the ordinary
analytic coefficient germs required by the global finite-halo theorem.


\section{Support-and-degree control beyond the finite halo}
\label{sec:domains}

\subsection{The rescaled input family}
Assume polynomial coefficients in \eqref{eq:Hinput}. Split each nonzero
coefficient into its finitely many monomials and regard a monomial occurrence
as an input label
\[
 \ell=(\gamma,\alpha,\beta),\qquad
 c_\ell t^\gamma q^\alpha p^\beta,
 \quad c_\ell\ne0,\quad d_\ell=|\alpha|+|\beta|\ge3.
\]
The label set is a set; infinitely many different labels are allowed. For
$\rho\in\Gamma$, define
\begin{equation}\label{eq:weights}
 w_\rho(\ell)=\gamma+(d_\ell-2)\rho,
 \qquad D_\rho=(t^\rho\OO_\Gamma)^{2d}.
\end{equation}

\begin{definition}[Admissible rescaling]\label{def:admissible}
The input monomial family is admissible at $\rho$ if
\begin{enumerate}[label=(\arabic*)]
\item $w_\rho(\ell)>0$ for every input label;
\item the set $W_\rho=\{w_\rho(\ell):\ell\}$ is well ordered;
\item each weight in $W_\rho$ comes from only finitely many input labels.
\end{enumerate}
These requirements assert strong summability of the individually rescaled
monomials into a positive Hahn series with polynomial coefficients.
\end{definition}
The point $\rho=0$ is always admissible. A negative $\rho$ gives a domain
containing phase points of infinite magnitude. For infinitely many input
monomials, condition (1) alone does not imply (2) or (3).

The Hamiltonian rescaling operator is
\begin{equation}\label{eq:S}
 \mathscr S_\rho F(z)=t^{-2\rho}F(t^\rho z).
\end{equation}
It is a Poisson \emph{Lie} algebra homomorphism wherever its monomial
expansion is summable:
\begin{equation}\label{eq:scalepb}
 \{\mathscr S_\rho F,\mathscr S_\rho G\}
       =\mathscr S_\rho\{F,G\}.
\end{equation}
It is not a homomorphism for ordinary products. Indeed, the factors
$t^{-2\rho}$ in \eqref{eq:S} are exactly what compensates for the two
phase derivatives in a Poisson bracket. The map fixes $H_0$, preserves
resonant monomials, and commutes with $L$ and $\Pi$.

\subsection{The domain theorem}
\begin{theorem}[Compatible support-certified domains]\label{thm:rescale}
Suppose the polynomial input of Theorem~\ref{thm:normalform} is admissible
at $\rho$. Then its canonical normalizing map and inverse have strongly
summable monomial realizations on $D_\rho$, and they are mutually inverse
symplectic bijections of that domain. The Hamiltonian identity and the
commuting action identities hold there.

If the normalizing construction is performed independently after rescaling
at two admissible values, its realized maps agree on the overlap after
undoing the scalings. Thus these domains carry one compatible normalization,
not a choice of unrelated local conjugacies.
\end{theorem}
\begin{proof}
Admissibility makes
\[
 H_\rho:=\mathscr S_\rho H
 =H_0+\sum_\ell c_\ell t^{w_\rho(\ell)}q^{\alpha_\ell}p^{\beta_\ell}
\]
a positive Hahn perturbation with polynomial coefficients. Apply
Theorem~\ref{thm:normalform}(a) to obtain $\chi_\rho$ and $\Phi_\rho$ on
the finite halo. Undoing the phase scaling gives a symplectic bijection
\begin{equation}\label{eq:Phirescaled}
 \Psi_\rho(z)=t^\rho\Phi_\rho(t^{-\rho}z)
 \quad\text{of }D_\rho.
\end{equation}
The equal scalar dilation in all phase coordinates is conformally
symplectic; the dilation and its inverse in \eqref{eq:Phirescaled} cancel,
so $\Psi_\rho$ is symplectic in the original bracket.

It remains to show that this is the same normalization, not merely some
map obtained from the rescaled problem. Use the labeled coefficient trees
from the proof of Theorem~\ref{thm:normalform}, now expanding each input
coefficient into its finitely many monomial labels. Every nonzero monomial
contribution to a generator or Hamiltonian coefficient with leaves
$\ell_1,\ldots,\ell_r$ has phase degree
\begin{equation}\label{eq:degreeham}
 2+\sum_{j=1}^r(d_{\ell_j}-2).
\end{equation}
Every contribution to a nonlinear coordinate coefficient has phase degree
\begin{equation}\label{eq:degreecoord}
 1+\sum_{j=1}^r(d_{\ell_j}-2).
\end{equation}
These formulas follow inductively because a Poisson bracket subtracts two
from the sum of phase degrees; $L$, $L^{-1}$, and $\Pi$ preserve degree.
The coordinate construction starts at degree one instead of two.

Under Hamiltonian scaling, the exponent of the tree in
\eqref{eq:degreeham} becomes precisely
\[
 \sum_{j=1}^r\gamma_{\ell_j}
    +\rho\sum_{j=1}^r(d_{\ell_j}-2)
 =\sum_{j=1}^r w_\rho(\ell_j).
\]
For coordinates the same sum is obtained by using
$t^{-\rho}\Phi(t^\rho z)$ instead. Since $W_\rho$ is positive and well
ordered, Lemma~\ref{lem:neumann} gives finitely many words at a fixed
rescaled weight. The finite-fiber hypothesis gives finitely many choices
of original labels for each word. There are then finitely many coefficient
trees. This proves the strong summability of the rescaled generator,
Hamiltonian, and coordinate expressions, before any regrouping is made.

Consequently \eqref{eq:scalepb} can be applied to every finite tree and then
summed. The rescaled generator is $\mathscr S_\rho\chi$; it satisfies the
same nonresonant gauge and normal-form equation as $\chi_\rho$. The
uniqueness theorem identifies them. On coordinates this says
\[
 \Phi_\rho(z)=t^{-\rho}\Phi(t^\rho z)
\]
where the right side is the supported monomial realization just proved.
The same argument applies to the inverse exponential and to the actions.
Taylor evaluation at finite rescaled points is valid by
Proposition~\ref{prop:eval}. This proves the identities on $D_\rho$.

For two admissible rescalings, both maps arise from the same labeled trees.
At a point of their overlap those tree evaluations are strongly summable
by each certificate, and every regrouping agrees coefficientwise. This
proves compatibility without assuming that rescaling preserves the order
of the original Hahn exponents; in general it does not.
\end{proof}

\begin{remark}[Scope of the extension]
The theorem extends the polynomial-coefficient normalizer, Hamiltonian,
and constructed actions. It does not automatically evaluate an arbitrary
ordinary entire function at an infinite phase or action argument. A
centralizer statement on $D_\rho$ can instead be made in the corresponding
\emph{rescaled} entire-coefficient algebra. Transporting all original
entire coefficients to infinite arguments without a summability certificate
would be an unjustified enlargement of the function class.
\end{remark}

\subsection{Homogeneous perturbations}
\begin{corollary}[A degree-dependent radius]\label{cor:homogeneous}
Let $H=H_0+t^\eta P_m$ with $\eta>0$ and $P_m$ a homogeneous polynomial of
phase degree $m\ge3$. Write the generator and coordinates by powers of
$t^\eta$:
\[
 \chi=\sum_{n\ge1}t^{n\eta}\chi_n,\qquad
 \Phi=z+\sum_{n\ge1}t^{n\eta}\Phi_n.
\]
Every nonzero $\chi_n$ is homogeneous of degree $2+n(m-2)$, and every
nonzero component of $\Phi_n$ has degree $1+n(m-2)$. The normalizer and
inverse are defined on every $D_\rho$ satisfying
\begin{equation}\label{eq:homogeneousdomain}
 \eta+(m-2)\rho>0.
\end{equation}
\end{corollary}
\begin{proof}
Every tree at exponent $n\eta$ has exactly $n$ leaves, because an ordered
abelian group is torsion-free and the only input weight is $\eta$. Apply
\eqref{eq:degreeham}--\eqref{eq:degreecoord}. The finitely many input
monomials all have the same shifted weight
$\eta+(m-2)\rho$, so its positivity gives all three admissibility
conditions.
\end{proof}
This is a guaranteed domain, not a general assertion that every point
outside it is singular. Special perturbations may admit larger or even
polynomial normalizers.

\subsection{A sharp cubic boundary and a global obstruction}
In one degree of freedom take
\begin{equation}\label{eq:cubicH}
 H(q,p)=qp+\varepsilon q^2p,\qquad \varepsilon=t^\eta,\quad\eta>0.
\end{equation}
There is no arithmetic issue because $\omega_1=1$. Put
$\chi=-\varepsilon q^2p$. Direct calculation gives
\begin{equation}\label{eq:cubicmap}
 \Phi(q,p)=\left(\frac{q}{1-\varepsilon q},\ p(1-\varepsilon q)^2\right),
 \qquad H\circ\Phi=qp.
\end{equation}
Indeed $\ad_\chi q=\varepsilon q^2$,
$\ad_\chi p=-2\varepsilon qp$, and the coordinate exponentials give
\eqref{eq:cubicmap}. One can also verify it by direct rational substitution.
Its Jacobian determinant is one, so it is symplectic. Since
$\Pi\chi=0$, this is the canonical normalizer of the theorem.

For $\eta+\rho>0$, the expansion
\[
 \frac{q}{1-\varepsilon q}=\sum_{n\ge0}\varepsilon^n q^{n+1}
\]
is strongly summable on $D_\rho$. At the boundary point
$q=\varepsilon^{-1}$, $p=0$, every term in this family has exponent
$-\eta$ and the rational function has a pole. Thus replacing the strict
inequality in \eqref{eq:homogeneousdomain} by a non-strict one cannot be
valid in general.

There is also a coordinate-independent obstruction on the whole boundary
polydisk $D_{-\eta}$. The Hamiltonian \eqref{eq:cubicH} has exactly two
critical points there:
\[
 (q,p)=(0,0),\qquad(q,p)=(-\varepsilon^{-1},0).
\]
This follows from $H_p=q(1+\varepsilon q)$ and
$H_q=p(1+2\varepsilon q)$. In contrast, $qp$ has exactly one critical
point. An everywhere invertible symplectic change on that domain satisfying
$H\circ\Psi=qp$ would give a bijection between these critical-point sets by
the chain rule, which is impossible. The boundary failure is therefore not
only a defect in one geometric-series proof.

\subsection{Two independent summability failures}
\begin{example}[Positive shifted weights, but no well-order]
\label{ex:badorder}
Over $\Gamma=\Q$, consider the valid polynomial-coefficient perturbation
\[
 R(q,p)=\sum_{n\ge2}t^{n+1/n}q^{n+2}.
\]
Its original exponents are strictly increasing. At $\rho=-1$ the shifted
weights are $1/n$, all strictly positive, but they form a descending
sequence. At the rescaled point $q=1$, the rescaled perturbation would be
$\sum_{n\ge2}t^{1/n}$, which is not a Hahn series. Positivity alone is
therefore insufficient even to define the rescaled input.
\end{example}

\begin{example}[Positive well-ordered weights, but infinite fibers]
\label{ex:badfiber}
Now take
\[
 R(q,p)=\sum_{n\ge2}t^{n+1}q^{n+2}.
\]
At $\rho=-1$ all shifted weights equal $1$. The weight set is positive and
well ordered, but infinitely many labels contribute to that one weight.
The rescaled expression is $t\sum_{n\ge2}q^{n+2}$, whose evaluation at the
ordinary point $q=1$ is not defined by finite coefficient addition or by
ordinary convergence. It cannot supply a polynomial-coefficient Hahn
perturbation on the whole rescaled finite halo.
\end{example}
These examples show why none of the three conditions in
Definition~\ref{def:admissible} can be dropped from a general theorem of
this form. They do not assert that finite fibers are necessary for every
possible analytic resummation in some larger, differently specified class.

\section{A genuinely multiscale worked example}
\label{sec:example}

\subsection{The Hamiltonian and its first coefficients}
Take $\omega=(1,\sqrt2)$, set $\kappa=2-\sqrt2$, and write
\[
 A=q_1^2p_2,\qquad B=p_1^2q_2,\qquad C=I_1I_2.
\]
For arbitrary positive scales $\eta,\theta\in\Gamma$, put
\begin{equation}\label{eq:workedH}
 H=I_1+\sqrt2 I_2+u(A+B)+vC,
 \qquad u=t^\eta,\quad v=t^\theta.
\end{equation}
Nonresonance follows from the irrationality of $\sqrt2$. The two scales need
not be Archimedean-comparable. For instance, the lexicographic example of
Remark~\ref{rem:rank} allows $n\eta<\theta$ for every $n$.

Temporarily retain $u$ and $v$ as independent \emph{labels}. The notation
$O_\lab(3)$ means terms involving at least three input labels, not terms
above a specified valuation cutoff. Direct homological calculation gives
\begin{align}
 \chi={}&\frac{u}{\kappa}(B-A)
 +\frac{uv}{\kappa^2}(A-B)(2I_2-I_1)+O_\lab(3),
 \label{eq:workedchi}\\
 N={}&I_1+\sqrt2 I_2+vI_1I_2
 +\frac{u^2}{\kappa}(I_1^2-4I_1I_2)+O_\lab(3).
 \label{eq:workedN}
\end{align}
The coefficients of $v$, $u^2$, and $v^2$ in $\chi$ vanish, and the
coefficients of $uv$ and $v^2$ in the displayed normal form vanish.

For clarity, the key bracket is
\[
 \{A,B\}=4I_1I_2-I_1^2.
\]
Since $LA=\kappa A$ and $LB=-\kappa B$, the first generator is
$u(B-A)/\kappa$. The second normal term is
\[
 \frac12\Pi\left\{\frac{u(B-A)}{\kappa},u(A+B)\right\}
 =\frac{u^2}{\kappa}(I_1^2-4I_1I_2).
\]
Also
\[
 \{A,C\}=A(2I_2-I_1),\qquad
 \{B,C\}=-B(2I_2-I_1),
\]
which gives the mixed generator in \eqref{eq:workedchi} by
$-L^{-1}Q$. These identities fix both the sign convention and the placement
of the inverse transformation.

\subsection{Coordinates, integrals, and certified domains}
To first label degree the coordinate map is
\begin{align*}
 \Phi(q_1)&=q_1-\frac{2u}{\kappa}p_1q_2+O_\lab(2),&
 \Phi(p_1)&=p_1-\frac{2u}{\kappa}q_1p_2+O_\lab(2),\\
 \Phi(q_2)&=q_2+\frac{u}{\kappa}q_1^2+O_\lab(2),&
 \Phi(p_2)&=p_2+\frac{u}{\kappa}p_1^2+O_\lab(2).
\end{align*}
The actions in original coordinates are
\begin{align*}
 J_1&=I_1+\frac{2u}{\kappa}(A+B)+O_\lab(2),\\
 J_2&=I_2-\frac{u}{\kappa}(A+B)+O_\lab(2).
\end{align*}
The full supported series satisfy $H=N(J)$ and commute exactly; the
remainders here are only for displaying a manageable initial segment.

There are two input degrees: the $u$ terms are cubic and the $v$ term is
quartic. Thus every generator or Hamiltonian coefficient with label $u^av^b$
is homogeneous of phase degree $2+a+2b$, while the corresponding coordinate
coefficient has degree $1+a+2b$. A sufficient domain certificate is simply
\begin{equation}\label{eq:exampledomain}
 \eta+\rho>0,\qquad \theta+2\rho>0.
\end{equation}
The input label set is finite, so positivity already supplies well-ordering
and finite fibers in this example. In a higher-rank group, writing these
as valuation inequalities rather than numerical radii avoids introducing a
nonexistent real-valued absolute value.

If two label weights happen to coincide, for example $2\eta=\theta$, their
terms are added at that Hahn exponent. The construction remains valid.
Label-degree truncation and Hahn-weight truncation are still different
operations; the verification code deliberately uses the former.

\section{A discrete one-variable parallel}
\label{sec:discrete}
The same support mechanism is visible in a simpler setting. We include it
to connect the Hamiltonian theory with established non-Archimedean
linearization, not as a claim that polynomial perturbative linearization is
new in rank one \cite{lindahl}.

\begin{proposition}[Polynomial-layer linearization]\label{prop:discrete}
Let $\lambda\in\C^\times$ be not a root of unity, and let
\[
 f(z)=\lambda z+R(z),\qquad R\in\HH{\Pp_1}_{>0},
 \quad\ord_0R_s\ge2.
\]
There is a unique $h(z)=z+U(z)$ with positive polynomial Hahn coefficients,
$\ord_0U_\gamma\ge2$, satisfying
\begin{equation}\label{eq:discreteconj}
 f\circ h=h\circ(\lambda\id).
\end{equation}
It is a bijection of the finite halo, and \eqref{eq:discreteconj} holds
there as an actual identity. No bound on $|\lambda^n-\lambda|^{-1}$ is
required.
\end{proposition}
\begin{proof}
Define $L_\lambda U=U(\lambda z)-\lambda U(z)$. On $z^n$, for $n\ge2$,
this operator is multiplication by the nonzero constant
$\lambda^n-\lambda$. It is therefore invertible on polynomials of order at
least two. Equation \eqref{eq:discreteconj} becomes
\[
 U=L_\lambda^{-1}R(z+U).
\]
At a fixed Hahn exponent, every nonlinear occurrence of $U$ on the right
is accompanied by a positive coefficient of $R$, hence uses strictly
smaller exponents of $U$. The supported recursion and finite-expression
argument of Proposition~\ref{prop:inverse} apply. Each coefficient remains
a polynomial, and the least-difference argument proves uniqueness.
Propositions~\ref{prop:inverse} and \ref{prop:eval} give the actual
conjugacy and bijection.
\end{proof}
For the example $f(z)=\lambda z+\varepsilon z^2$, writing
\[
 h(z)=\sum_{n\ge1}b_n\varepsilon^{n-1}z^n,\qquad b_1=1,
\]
gives the exact recurrence
\begin{equation}\label{eq:discreterec}
 (\lambda^n-\lambda)b_n=\sum_{j=1}^{n-1}b_jb_{n-j}\qquad(n\ge2).
\end{equation}
There is one phase monomial per Hahn layer. At finite $z$, the support
rather than the complex growth of the $b_n$ is what makes this a valid Hahn
identity. The supplied verifier checks \eqref{eq:discreterec} exactly for
$\lambda=2$ through $n=10$. This finite check says nothing about all $n$;
the coefficient derivation above proves the recurrence.

\section{Reproducibility, proof audit, and further questions}
\label{sec:verification}

\subsection{What the executable checks do}
The accompanying \texttt{code/verify.py} works over the exact field
$\Q(\sqrt2)$ and truncates the independent label ring $\Q(\sqrt2)[[u,v]]$
at total label degree four. It constructs the single-generator normal form
for \eqref{eq:workedH} and checks the following kinds of identities:
\begin{enumerate}[label=(\roman*)]
\item every computed normal coefficient is resonant and every generator
coefficient satisfies the nonresonant gauge;
\item the displayed first and mixed coefficients, symplectic coordinate
brackets, inverse Lie maps, direct substitution $H(\Phi)=N$, and
$J(\Phi)=I$ agree exactly;
\item the actions commute with the Hamiltonian and one another, a nonlinear
function of them is a first integral, and all computed coefficients obey
the degree-defect rules;
\item the rational cubic boundary map preserves the Hamiltonian and
symplectic form, the support counterexample has the claimed finite-sample
order pattern, and the discrete recurrence holds at the tested degrees.
\end{enumerate}
The recorded execution is
\begin{quote}\ttfamily
PASS: 119 exact checks through total label degree 4.
\end{quote}
All arithmetic identities use exact symbolic operations, not numerical
tolerances. The JSON record includes every check name and the computed
generator and normal-form coefficients.

The script is not a formal proof of arbitrary-support summability, the
entire-coefficient threshold, the centralizer theorem, or the Liouville
construction. Those are established by the arguments in this article, not
by testing. No Lean or other proof-assistant formalization is supplied or
claimed.

\subsection{A dependency and scope audit}
The logical dependence is deliberately short. Neumann's positive-support
lemma supplies every Hahn summation and reordering. Ordinary Taylor and
Cauchy theory supplies finite-halo evaluation and the entire action
projection. The sharp divisor estimate says exactly when the homological
inverse remains in the entire coefficient algebra. The normal-form fixed
point then uses only those operations. The centralizer theorem is a
least-exponent argument after that normalization. The larger-domain theorem
adds a separate monomial support certificate; it is not an automatic
extension of ``entire'' to infinite arguments.

The following boundaries are essential to interpreting the results.
\begin{itemize}
\item No ordinary analytic convergence in a complex perturbation parameter
is inferred. A substitution $t^\eta=\varepsilon\in\C^\times$ need not exist
on the coefficient algebra and may turn a valid Hahn series into a
divergent ordinary series.
\item No general statement is made for unrelated germ coefficients, for
infinite-dimensional Hamiltonian systems, or for resonant frequencies.
The gauge equation changes when its homological inverse has a larger
kernel.
\item No fixed real-valued norm or sequential completeness is assumed of
$K_\Gamma$. Arbitrary rank is handled by supports. The full class of
surcomplex numbers is used only through compatible set-sized workspaces.
\end{itemize}
These are not technical qualifications appended to a stronger theorem; they
are part of the mathematical category in which the theorem is true.

\subsection{Further questions suggested by the proved results}
Several concrete problems remain beyond the statements established here.
For resonant frequency vectors, the first reduction retains a nontrivial
resonant Hamiltonian algebra. It is natural to ask for effective support
conditions under which the dynamics of that retained algebra can itself be
integrated, but the present nonresonant centralizer proof does not answer
that question.

For entire layers failing the universal threshold, particular inputs may
still normalize because their coefficients vanish at the dangerous
monomials. An input-specific necessary and sufficient condition would have
to remain stable under the nonlinear brackets created by normalization;
a first-order divisibility test alone need not be enough.

Finally, the rescaling theorem gives guaranteed domains, not maximal ones.
The cubic example detects a sharp boundary through both a pole and an
extra critical point. A general procedure combining support certificates
with algebraic or analytic obstructions could yield substantially better
maximal-domain information for finitely presented perturbations.

\subsection{Conclusion}
Hahn summability does not merely reproduce a divergent formal normal form
under a new name. For polynomial layers, and for entire layers precisely
under the universal exponential-divisor condition, it realizes the normal
form as an invertible analytic map on all finite surcomplex phase points.
The realization supplies exact commuting actions and an exhaustive
entire-coherent centralizer. At infinite scales, the correct additional
hypothesis is a joint support-and-degree condition, not positivity of each
term in isolation. The counterexamples identify where each of these
conclusions stops being valid.

\appendix
\section{Coefficient recursion as a reproducible algorithm}
\label{app:algorithm}
For a finite collection of input labels $x_1,\ldots,x_r$, an implementation
can temporarily work over the ordinary formal label ring. The following
algorithm computes through a chosen total label degree $N$.
\begin{enumerate}[label=\arabic*.]
\item Set $\chi=0$ and retain $H=H_0+R$ with independent labels attached to
its input layers.
\item For $n=1,\ldots,N$, compute $e^{\ad_\chi}H$ through total label degree
$n$, using the current $\chi$ of label degree less than $n$.
\item For each coefficient $E_\nu$ with $|\nu|=n$, append
$-L^{-1}QE_\nu$ to the coefficient of $x^\nu$ in $\chi$.
\item Recompute $N_H=e^{\ad_\chi}H$ through the requested order, and compute
coordinates by the same exponential and actions by the inverse exponential.
\end{enumerate}
Adding a degree-$n$ generator changes the degree-$n$ transformed Hamiltonian
only by $L\chi_n$. Thus step 3 is the single-exponential gauge recursion,
not the different algorithm that composes a list of successive exponential
maps without taking their logarithm.

For an infinite input label set, this is a conceptual construction rather
than a finite executable loop over all inputs. Homogeneous leaf degree can
still organize the universal expression trees. At a specified Hahn weight,
Neumann's lemma permits only finitely many leaf words, hence only finitely
many homogeneous leaf degrees and finite expressions. That is the exact
connection between the finite-label algorithm and the arbitrary-support
proof.

In the accompanying two-label test, total label degree four is not a
valuation cutoff. If $\theta>n\eta$ for every $n$, infinitely many powers of
$u$ lie below the first $v$ layer, and no finite label truncation pretends to
compute that entire valuation initial segment.

\section{Build and verification instructions}
The archive contains a standalone \texttt{article.tex} with an internal
bibliography, its compiled \texttt{article.pdf}, the exact verifier, its JSON
record, and a README. No external figures, bibliography database, private
repository access, or proprietary computer algebra system are required.

From the archive directory, compile with
\begin{verbatim}
latexmk -pdf -interaction=nonstopmode -halt-on-error article.tex
\end{verbatim}
Alternatively run \texttt{pdflatex} repeatedly until references stabilize.
The source uses common packages from a standard TeX Live installation.

With Python 3.10 or later and SymPy installed, run
\begin{verbatim}
python code/verify.py
\end{verbatim}
The verifier raises an error on a failed identity and rewrites
\texttt{data/verification.json} only after all checks pass. The supplied
run used SymPy 1.14.0. Reproduction of the finite checks is not certification
of the infinite mathematical statements.

\begin{thebibliography}{9}
\raggedright
\addcontentsline{toc}{section}{References}

\bibitem{repository}
V. Reshetnikov, \emph{Surreal: surreal and surcomplex research reports},
repository catalogue and foundational analysis README, AI-assisted drafts,
commit \texttt{39f2be6667ade51bca2b45daa47e289d69c09764}, inspected
21 September 2026.
\url{https://github.com/VladimirReshetnikov/Surreal/tree/39f2be6667ade51bca2b45daa47e289d69c09764/docs}.
The present comparison uses the catalogue and analysis README, not an
assertion that every source manuscript has been exhaustively audited.

\bibitem{neumann}
B. H. Neumann, \emph{On ordered division rings},
Transactions of the American Mathematical Society \textbf{66} (1949),
202--252.
\url{https://doi.org/10.1090/S0002-9947-1949-0032593-5}.

\bibitem{perezmarco}
R. P\'erez-Marco, \emph{Convergence or generic divergence of Birkhoff normal
form}, arXiv:math/0009028 (2000).
\url{https://arxiv.org/abs/math/0009028}.

\bibitem{krikorian}
R. Krikorian, \emph{On the divergence of Birkhoff Normal Forms},
arXiv:1906.01096 (2019).
\url{https://arxiv.org/abs/1906.01096}.

\bibitem{lindahl}
K.-O. Lindahl, \emph{Linearization in ultrametric dynamics in fields of
characteristic zero---equal characteristic case},
p-Adic Numbers, Ultrametric Analysis and Applications \textbf{1}(4)
(2009), 307--316; arXiv:1111.1993.
\url{https://arxiv.org/abs/1111.1993}.

\bibitem{bagayoko}
V. Bagayoko, L. S. Krapp, S. Kuhlmann, D. Panazzolo, and M. Serra,
\emph{Automorphisms and derivations on algebras endowed with formal infinite
sums}, arXiv:2403.05827v2 (2025 revision).
\url{https://arxiv.org/abs/2403.05827v2}.

\bibitem{kaplan}
E. Kaplan, L. S. Krapp, and M. Serra,
\emph{Decomposing the automorphism group of the surreal numbers},
arXiv:2509.22374v3 (2026 revision).
\url{https://arxiv.org/abs/2509.22374v3}.

\end{thebibliography}
\end{document}
